\DeclareMathOperator{\ones}{\mathbf 1}
\DeclareMathOperator{\R}{\mathbb{R}}
\DeclareMathOperator{\C}{\mathbb{C}}
\DeclareMathOperator{\Z}{\mathbb{Z}}
\DeclareMathOperator{\N}{\mathbb{N}}
\DeclareMathOperator{\Q}{\mathbb{Q}}
\DeclareMathOperator{\A}{\mathbb{A}}
\DeclareMathOperator{\symm}{{\mbox{\bf S}}}  
\DeclareMathOperator{\diam}{{\mbox{diam}}}  
\DeclareMathOperator{\Conv}{{\mbox{Conv}}}  
\DeclareMathOperator{\SLE}{{\mbox{SLE}}}  

\DeclareMathOperator{\nullspace}{{\mathcal N}}
\DeclareMathOperator{\range}{{\mathcal R}}
\DeclareMathOperator{\Rank}{\mathop{\bf Rank}}
\DeclareMathOperator{\Tr}{\mathop{\bf Tr}}
\DeclareMathOperator{\diag}{\mathop{\bf diag}}
\DeclareMathOperator{\card}{\mathop{\bf card}}
\DeclareMathOperator{\rank}{\mathop{\bf rank}}
\DeclareMathOperator{\conv}{\mathop{\bf conv}}
\DeclareMathOperator{\prox}{\mathbf{prox}}
\DeclareMathOperator{\per}{\mbox{per}}



\DeclareMathOperator{\Var}{\mathop{\bf Var{}}}
\DeclareMathOperator{\Cov}{\mathop{\bf Cov{}}}
\DeclareMathOperator{\Prob}{\mathop{\bf Prob}}
\DeclareMathOperator{\Co}{{\mathop {\bf Co}}} % convex hull
\DeclareMathOperator{\dist}{\mathop{\bf dist{}}}
\DeclareMathOperator{\argmin}{\mathop{\rm argmin}}
\DeclareMathOperator{\argmax}{\mathop{\rm argmax}}
\DeclareMathOperator{\epi}{\mathop{\bf epi}} % epigraph
\DeclareMathOperator{\Vol}{\mathop{\bf vol}}
\DeclareMathOperator{\dom}{\mathop{\bf dom}} % domain
\DeclareMathOperator{\intr}{\mathop{\bf int}}
\DeclareMathOperator{\sign}{\mathop{\bf sign}}

\newcommand{\cf}{{\it cf.}}
\newcommand{\eg}{{\it e.g.}}
\newcommand{\ie}{{\it i.e.}}
\newcommand{\etc}{{\it etc.}}
\newcommand{\etal}{\textit{et al}. }
\DeclareMathOperator{\len}{\mathop{len}}
%%additional commands



%%theorems 
\newtheorem{theorem}{Theorem}[section]
\newtheorem{prop}[theorem]{Proposition}
\newtheorem{fact}[theorem]{Fact}
\newtheorem{cor}[theorem]{Corollary}
\newtheorem{lemma}[theorem]{Lemma}
\newtheorem{claim}{Claim}
\newtheorem{remark}[theorem]{Remark}
\newtheorem{definition}[theorem]{Definition}
\newtheorem{problem}[theorem]{Problem}
\newtheorem{conj}[theorem]{Conjecture}

%% interior and closure
\newcommand{\cl}{\mathbf{cl}}
\DeclareMathOperator{\inte}{\mathbf{int}}

%% redefine P to refer to probability
\renewcommand*{\P}{\mathbb P}
\DeclareMathOperator{\E}{\mathop{\mathbb E{}}}


\renewcommand*{\S}{\mathbb S}
\DeclareMathOperator*{\osc}{osc}


%to stop the bib warning
\renewcommand{\Bbb}{\mathbb}

\newcommand{\indc}{{\boldsymbol{1}}}
\newcommand*{\Zd}{\ensuremath{\mathbb{Z}^d}}
\newcommand*{\Rd}{\ensuremath{\mathbb{R}^d}}
\newcommand*{\supp}{\ensuremath{\mathrm{supp\,}}}
\newcommand{\cu}{\square}
\newcommand{\size}{\mathrm{size}}

%% homogenization preliminaries
\renewcommand{\a}{\mathbf{a}}
\newcommand{\Abf}{{\mathbf{A}}}

\newcommand*\xxbar[1]{%
	\hbox{%
		\vbox{%
			\hrule height 0.7pt % The actual bar
			\kern0.25ex%         % Distance between bar and symbol
			\hbox{%
				\kern+0em%      % Shortening on the left side
				\ensuremath{#1}%
				\kern-0.1em%      % Shortening on the right side
			}%
		}%
	}%
}
\newcommand{\ahom}{\xxbar{\mathbf{a}}\,}



\newcommand{\Abfh}{{\overbracket[.8pt][-1pt]{{\mathbf{A}}}}}  % change the first number to change the size of the bar
\newcommand{\Bhom}{{\overbracket[.8pt][-1pt]{{\mathbf{B}}}}}  % change the first number to change the size of the bar


